\par\medskip\noindent
\underline{Lema 2.15.} Sean $A$ y $B$ grupoides discretos sobre
$\underline C_\Lambda$, y sea $\varphi:A\longrightarrow B$ un funtor
mínimamente liso. Si $A$ y $B$ son homogéneos, entonces $\varphi$ es una
$\underline C_\Lambda$-equivalencia.

\underline{Demostración.} Debemos demostrar que
$\varphi(S):A(S)\longrightarrow B(S)$ es biyectiva para todo
$S\in\operatorname{ob}(\underline C_\Lambda)$. Es suprayectiva para todo
$S$ y biyectiva para $S=k[\varepsilon]$. Razonamos por inducción sobre la
longitud de $S$. Sea $S'\xrightarrow fS$ una extensión pequeña en
$\underline C_\Lambda$. Puesto que $A$ y $B$ son grupoides discretos
homogéneos, el isomorfismo canónico
\[
 1\times\widetilde f:S'\times_SS'
 \xrightarrow{\ \sim\ }S'\times_k k[\operatorname{Ker}f]
\]
induce un diagrama conmutativo

%%%%%%%%%%  scan idx 78  |  folio 67  %%%%%%%%%%

\[
\begin{tikzcd}[column sep=huge,row sep=huge]
 A(S')\times_{A(S)}A(S') \arrow[r,"\sim"] \arrow[d]
 & A(S')\times A(k[\operatorname{Ker}f]) \arrow[d]\\
 B(S')\times_{B(S)}B(S') \arrow[r,"\sim"']
 & B(S')\times B(k[\operatorname{Ker}f]).
\end{tikzcd}
\]
Así, $A(S')$ y $B(S')$ son conjuntos homogéneos principales sobre
$A(S)$ y $B(S)$, con grupos estructurales identificados por
\[
 A(k[\operatorname{Ker}f])\xrightarrow{\ \sim\ }
 B(k[\operatorname{Ker}f]).
\]
Si $\varphi(S)$ es biyectiva, se sigue que $\varphi(S')$ es inyectiva y,
por tanto, biyectiva. Esto concluye la inducción.
\hfill$\square$

\underline{Corolario 2.16.} Sea $\underline A$ un grupoide homogéneo
sobre $\underline C_\Lambda$ tal que
\[
 \operatorname{Aut}_{\underline A(k[\varepsilon])}(1_\varepsilon)
\]
sea de dimensión finita. Para todo
$\xi\in\operatorname{ob}(\widehat{\underline A})$, el funtor
\[
 \operatorname{Isom}_\xi:
 \underline C_{R\widehat\otimes_\Lambda R}
 \longrightarrow(\mathrm{Sets})
\]
es prorepresentable, donde $R=P(\xi)$.

\underline{Demostración.} Si $\operatorname{Isom}_\xi$ es homogéneo,
el Teorema 1.11 proporciona un funtor mínimamente liso
\[
 \Phi:\underline h_{R_1}\longrightarrow
 \underline{\operatorname{Isom}}_\xi.
\]
El Lema 2.15 hace de $\Phi$ una
$\underline C_{R\widehat\otimes_\Lambda R}$-equivalencia,
y por tanto $\operatorname{Isom}_\xi$ es prorepresentable. Queda por
demostrar la homogeneidad. Considérese un diagrama
\[
\begin{tikzcd}[column sep=large,row sep=large]
 T_1 \arrow[dr] && T_2 \arrow[dl]\\
 & T &
\end{tikzcd}
\]
en $\underline C_{R\widehat\otimes_\Lambda R}$, con
$T_2\longrightarrow T$ suprayectiva. Puesto que $\underline A$ es
homogéneo, los productos fibrados
\[
 \xi_i(T_1)\times_{\xi_i(T)}\xi_i(T_2)
\]
existen en $\underline A$, y hay isomorfismos canónicos
\[
 \xi_i(T_1\times_TT_2)\xrightarrow{\ \sim\ }
 \xi_i(T_1)\times_{\xi_i(T)}\xi_i(T_2)
 \qquad(i=1,2)
\]
(cf. 2.6(a)). Por consiguiente,
\[
 \operatorname{Isom}_\xi(T_1\times_TT_2)
 \longrightarrow
 \operatorname{Isom}_\xi(T_1)
 \times_{\operatorname{Isom}_\xi(T)}
 \operatorname{Isom}_\xi(T_2)
\]
es un isomorfismo. \hfill$\square$

\underline{Teorema 2.17.} Sea $\underline A$ un grupoide cofibrado
sobre $\underline C_\Lambda$. Entonces es prorepresentable por un
cogrupoide liso (normalizado) si y solo si

%%%%%%%%%%  scan idx 79  |  folio 68  %%%%%%%%%%

(1) $\underline A$ es homogéneo;

(2) $[\underline A(k[\varepsilon])]$ y
$\operatorname{Aut}_{\underline A(k[\varepsilon])}(1_\varepsilon)$
son espacios vectoriales de dimensión finita sobre $k$.

\underline{Demostración.} Supongamos primero que $\underline A$ está
prorepresentado por un cogrupoide liso. Así, $\underline A$ es
$\underline C_\Lambda$-equivalente a
\[
 \underline h=\underline h_{(R_0,R_1,\ldots)}
\]
para un cogrupoide liso $(R_0,R_1,\ldots)$ en
$\widehat{\underline C}_\Lambda$. La Proposición 2.12(a) muestra que
$\underline h$ es homogéneo. Además,
\begin{align*}
 [\underline h(k[\varepsilon])]
 &=\operatorname{Coker}\!\left(
   \operatorname{Hom}_{\Lambda\text{-alg}}(R_1,k[\varepsilon])
   \mathrel{\substack{\xrightarrow{\ s_1\ }\\[-.6ex]
                       \xrightarrow[\ s_2\ ]{}}}
   \operatorname{Hom}_{\Lambda\text{-alg}}(R_0,k[\varepsilon])
   \right),\\[1ex]
 \operatorname{Aut}_{\underline h}(1_\varepsilon)
 &=\operatorname{Ker}\!\left(
   \operatorname{Hom}_{\Lambda\text{-alg}}(R_1,k[\varepsilon])
   \xrightarrow{\ s_1\ }
   \operatorname{Hom}_{\Lambda\text{-alg}}(R_0,k[\varepsilon])
   \right)\\
 &=\operatorname{Ker}\!\left(
   \operatorname{Hom}_{\Lambda\text{-alg}}(R_1,k[\varepsilon])
   \xrightarrow{\ s_2\ }
   \operatorname{Hom}_{\Lambda\text{-alg}}(R_0,k[\varepsilon])
   \right).
\end{align*}
Ambos son espacios vectoriales de dimensión finita sobre $k$, por lo que
lo mismo vale para $\underline A$.

Recíprocamente, supongamos (1) y (2). Como $\underline A$ es homogéneo y
$[\underline A(k[\varepsilon])]$ es de dimensión finita, posee un objeto
formal mínimamente versal $\xi$. Pongamos $P(\xi)=R_0$ y consideremos
\[
 \operatorname{Isom}_\xi:
 \widehat{\underline C}_{R_0\widehat\otimes_\Lambda R_0}
 \longrightarrow(\mathrm{Sets}),
\]
dado por
\[
 T\longmapsto
 \operatorname{Isom}_{\underline A(T)}(\xi_1(T),\xi_2(T)),
\]
donde $\xi_i(T)=\xi(f_i)$ y $f_i:R_0\longrightarrow T$ son las
aplicaciones compuestas en
\[
\begin{tikzcd}[column sep=large]
 R_0 \arrow[r,shift left=2,"s_1"]
     \arrow[r,shift right=2,"s_2"']
 & R_0\widehat\otimes_\Lambda R_0 \arrow[r] & T.
\end{tikzcd}
\]
El Corolario 2.16 prorepresenta este funtor mediante un objeto
\[
 R_1\in
 \operatorname{ob}(\widehat{\underline C}_{R_0\widehat\otimes_\Lambda R_0})
\]
y un isomorfismo
$\xi_1(R_1)\xrightarrow{\varphi}\xi_2(R_1)$ para el cual
\[
 \Phi:\underline h_{R_1}\longrightarrow
 \underline{\operatorname{Isom}}_\xi
\]
es un isomorfismo sobre
$\widehat{\underline C}_{R_0\widehat\otimes_\Lambda R_0}$. La pareja
$(R_0,R_1)$ define, por tanto, un cogrupoide, y
\[
 (\xi,\varphi):\underline h_{(R_0,R_1,\ldots)}
 \longrightarrow\underline A
\]
es una $\underline C_\Lambda$-equivalencia. Como $\xi$ es mínimamente
versal, el compuesto
\[
 \underline h_{(R_0,R_0,\ldots)}
 \xrightarrow{\ \epsilon\ }
 \underline h_{(R_0,R_1,\ldots)}
 \xrightarrow{\ (\xi,\varphi)\ }
 \underline A
\]
es mínimamente liso. Puesto que $(\xi,\varphi)$ es una equivalencia,
$\epsilon$ es asimismo mínimamente liso. Así,
$(R_0,R_1,\ldots)$ es un cogrupoide liso normalizado en
$\widehat{\underline C}_\Lambda$. \hfill$\square$
